% --------------------------------------------------------------
% This part is the preamble, which you don't have to worry about
% To type your solutions, scroll down to where it says "Start here"
% --------------------------------------------------------------

% This defines the formatting for the document
\documentclass[12pt]{article}
\usepackage[left=1in,top=1in,right=1in,bottom=1in]{geometry}
\usepackage{enumitem}
\setlist{noitemsep}
\setlist[enumerate,1]{label=(\alph*)}
\setlist[enumerate,2]{label=(\roman*)}

% This imports the packages which we will need to type math symbols
\usepackage{amsmath,amssymb}

% This defines the "problem" and "solution" environments
\usepackage{amsthm}
\theoremstyle{definition}\newtheorem{problem}{Problem}
\newenvironment{solution}{\begin{proof}[\bfseries\textup{Solution:}]}{\end{proof}}

\begin{document}

% --------------------------------------------------------------
%                         Start here
% --------------------------------------------------------------

% This is the title:
\begin{center}
\bfseries Math 521, Section 2
\\ 
Homework 1
\\ 
(Name) % Fill in your name here
\\ [24pt] 
\end{center}

% Problem 1 -----------------------------------------------------
\begin{problem}
Using induction, prove that for all $n\in\mathbb{N}$ we have
\begin{equation*}
3+11+\dots +(8n-5)=4n^2-n.
\end{equation*}
\end{problem}
\begin{solution}
(Type your solution to problem 1 here.)
\end{solution}

% Problem 2 -----------------------------------------------------
\newpage
\begin{problem}
Compute the following sum:
\begin{equation*}
\frac{1}{1\cdot 4} + \frac{1}{4\cdot 7} + \dots + \frac{1}{(3n-2)(3n+1)} .
\end{equation*}
Prove that your answer is true for all $n\in\mathbb{N}$ using induction.
\end{problem}
\begin{solution}
(Type your solution to problem 2 here.)
\end{solution}

% Problem 3 -----------------------------------------------------
\newpage
\begin{problem}
Let $f:A\to B$ be a function.  
\begin{enumerate}
\item Prove that $f$ is injective if and only if $f^{-1}(f(X)) = X$ for all $X\subseteq A$.
\item Prove that $f$ is surjective if and only if $f(f^{-1}(Y)) = Y$ for all $Y\subseteq B$.
\end{enumerate}
\end{problem}
\begin{solution}
(a) (Type your solution to problem 3a here.)

(b) (Type your solution to problem 3b here.)

(You may want to use the following symbols: `for all' $\forall$, `there exists' $\exists$, and `in' $\in$.)
% The dollar signs "$" start and end inline math mode; these symbols can only be used in math mode. 
\end{solution}

% Problem 4 -----------------------------------------------------
\newpage
\begin{problem}
Let $f : A\to B$ be a function.  Prove that the following two statements are equivalent:
\begin{enumerate}
\item The function $f$ is surjective.
\item For any set $C$ and for any functions $g: B \to C$ and $h : B \to C$ such that $g\circ f = h\circ f$, we have $g = h$.
\end{enumerate}
\end{problem}
\begin{solution}
(Type your solution to problem 4 here.)
\end{solution}

% Problem 5 -----------------------------------------------------
\newpage
\begin{problem}
Let $f : B\to C$ be a function.  Prove that the following two statements are equivalent:
\begin{enumerate}
\item The function $f$ is injective.
\item For any set $A$ and for any functions $g: A \to B$ and $h : A \to B$ such that $f\circ g = f\circ h$, we have $g = h$.
\end{enumerate}
\end{problem}
\begin{solution}
(Type your solution to problem 5 here.)
\end{solution}

% Problem 6 -----------------------------------------------------
\newpage
\begin{problem}
Let $f:A\to B$ be a function.  Suppose that there exists a function $g:B\to A$ such that $(g\circ f)(a) = a$ for all $a\in A$ and $(f\circ g)(b) = b$ for all $b\in B$.  Prove that $f$ is bijective.
\end{problem}
\begin{solution}
(Type your solution to problem 6 here.)
\end{solution}

\end{document}